% ============================================================
%  CHƯƠNG IV: KẾT LUẬN VÀ HƯỚNG PHÁT TRIỂN
% ============================================================
\chapter{KẾT LUẬN VÀ HƯỚNG PHÁT TRIỂN}
\label{chap:ket_luan}

% ------------------------------------------------------------
\section{Những kết quả đã đạt được}
\label{sec:ket_qua_dat_duoc}

Sau thời gian nghiên cứu và phát triển, ứng dụng CineFlow đã cơ bản hoàn thiện và đáp ứng được các mục tiêu cốt lõi đặt ra ở Chương I cũng như tuân thủ các nguyên tắc thiết kế đã trình bày ở Chương II. Bảng~\ref{tab:ket_qua_doi_chieu} dưới đây đối chiếu trực tiếp giữa mục tiêu ban đầu và kết quả thực tế đạt được:

\begin{table}[H]
    \centering
    \renewcommand{\arraystretch}{1.5}
    \caption{Bảng đối chiếu mục tiêu và kết quả thực hiện}
    \label{tab:ket_qua_doi_chieu}
    \begin{tabularx}{\textwidth}{|>{\bfseries\raggedright\arraybackslash}p{4cm}|X|>{\centering\arraybackslash}p{2cm}|}
        \hline
        \rowcolor{gray!30} \textbf{Mục tiêu ban đầu} & \textbf{Kết quả đạt được} & \textbf{Tỉ lệ (\%)} \\ \hline
        Xây dựng ứng dụng xem phim trực tuyến trên Android & Hoàn thiện luồng từ Đăng ký/Đăng nhập, Trang chủ, Chi tiết phim, Phát video đến Shorts và lịch Premier League. & 100\% \\ \hline
        Áp dụng kiến trúc Client--Server phân lớp & Backend tổ chức rõ rệt ba tầng Controller--Service--Repository; Android client phân lớp UI--Repository--Network. & 100\% \\ \hline
        Quản lý nội dung tập trung qua Admin Panel & Triển khai đầy đủ CRUD phim/danh mục/người dùng và thống kê dữ liệu trực quan qua \texttt{AdminDashboardActivity}, \texttt{AdminFilmListActivity}, \texttt{AdminCategoryListActivity}, \texttt{AdminUserListActivity} và \texttt{AdminAnalyticsActivity}. & 100\% \\ \hline
        Phát video streaming mượt mà & ExoPlayer (Media3 1.5) streaming trực tiếp từ URL nguồn tĩnh trong MinIO Object Storage hoặc liên kết ngoài, không qua backend. & 100\% \\ \hline
        Bảo mật phiên đăng nhập & JWT Bearer Token (2 giờ) kết hợp cơ chế Refresh Token xoay vòng (Rotate Token) bảo mật, mật khẩu băm BCrypt, tự động đính kèm token qua Interceptor. & 100\% \\ \hline
    \end{tabularx}
\end{table}

Bên cạnh việc đáp ứng các yêu cầu nghiệp vụ, đồ án đã chứng minh được năng lực vận dụng nhiều mẫu thiết kế và công nghệ tiêu biểu trong lập trình di động:
\begin{itemize}
    \item \textbf{Kiến trúc phân lớp ba tầng} ở backend (Spring Boot 3.5.11) được tách bạch rõ ràng giữa Controller, Service và Repository, kết hợp \textbf{DTO Pattern} và vỏ phản hồi thống nhất \texttt{ApiResponse<T>}.
    \item \textbf{Quản lý schema bằng Flyway} với các phiên bản di trú (từ \texttt{V1} đến \texttt{V17}), đảm bảo schema PostgreSQL luôn nhất quán giữa máy của các thành viên và môi trường demo.
    \item \textbf{Spring Security Filter Chain} kết hợp \texttt{AuthTokenFilter} và \texttt{JwtEntryPoint} cho cơ chế xác thực JWT stateless, bảo vệ phân quyền chặt chẽ ở tầng URL (\texttt{/admin/**} chỉ cho phép \texttt{ROLE\_ADMIN}, \texttt{/user/**} yêu cầu xác thực).
    \item Phía Android áp dụng đồng thời nhiều mẫu thiết kế: \textbf{Repository Pattern}, \textbf{Singleton} (\texttt{ApiClient}, \texttt{AuthManager}), \textbf{Interceptor} (\texttt{AuthInterceptor}), \textbf{Observer} qua \texttt{LiveData} và \textbf{Template Method} qua \texttt{BaseFragment}.
    \item \textbf{Phát video tách biệt khỏi luồng dữ liệu nghiệp vụ}: ExoPlayer streaming trực tiếp từ URL nguồn (HTTP/HLS) lưu trữ tại MinIO hoặc liên kết ngoài, backend chỉ cung cấp \texttt{videoUrl} trong DTO -- giảm đáng kể tải cho backend.
\end{itemize}

\textbf{Kết luận chung:} Ứng dụng CineFlow đã giải quyết thành công bài toán cốt lõi của một nền tảng xem phim trực tuyến trên thiết bị Android, đồng thời cung cấp đầy đủ công cụ quản trị nội dung cho Admin. Đồ án không chỉ dừng lại ở mức một ứng dụng CRUD đơn giản mà đã hiện thực hoá đầy đủ luồng xác thực, quản lý nội dung, phát video và các mẫu thiết kế chuẩn mực của lập trình di động hiện đại.

% ------------------------------------------------------------
\section{Hạn chế và hướng phát triển}
\label{sec:han_che_huong_phat_trien}

Mặc dù đã đạt được những kết quả nhất định, do giới hạn về thời gian và phạm vi của một đồ án môn học, hệ thống vẫn còn một số hạn chế kỹ thuật:

\subsection{Hạn chế hiện tại}
\begin{itemize}
    \item \textbf{Tích hợp cổng thanh toán thực tế:} Hệ thống đã có cơ chế quản lý gói dịch vụ và kiểm tra quyền hạn thuê bao (\texttt{isPremium}), tuy nhiên chưa kết nối với các cổng thanh toán trực tuyến thực tế (như VNPay, MoMo) để kích hoạt tự động mà mới đang ở dạng giả lập (mock).
    \item \textbf{Phân phối nội dung video:} Hệ thống truyền tải video trực tiếp qua các file định dạng tĩnh (MP4) từ hệ thống lưu trữ MinIO hoặc URL ngoài, chưa áp dụng cơ chế truyền tải động thích ứng phân đoạn (như HLS, DASH) và dịch vụ CDN để tối ưu hóa băng thông theo thiết bị.
    \item \textbf{Môi trường triển khai:} Backend mới chạy ở \texttt{localhost:8080} (HTTP plaintext) phục vụ emulator qua \texttt{10.0.2.2}; chưa có cấu hình HTTPS, reverse proxy và đóng gói triển khai cho thiết bị thật trên LAN/Internet.
    \item \textbf{Kiểm thử tự động:} Đồ án chưa xây dựng bộ unit test/integration test cho backend và Android client; toàn bộ việc kiểm tra hiện được thực hiện thủ công qua Postman và emulator.
\end{itemize}

\subsection{Hướng phát triển ngắn hạn (3--6 tháng)}
Để củng cố chất lượng sản phẩm và đưa hệ thống tiệm cận mức triển khai thử nghiệm thực tế, nhóm đề ra các hướng phát triển tiếp theo:
\begin{itemize}
    \item \textbf{Tích hợp cổng thanh toán trực tuyến:} Kết nối API các gói thuê bao (\texttt{SubscriptionPackage}) với các dịch vụ thanh toán phổ biến như VNPay hoặc MoMo để tự động hóa quy trình mua và kích hoạt tài khoản Premium của người dùng.
    \item \textbf{Tối ưu hóa truyền tải video thích ứng:} Sử dụng giao thức phát video phân đoạn thích ứng (HLS/DASH) tích hợp vào ExoPlayer để tự động thay đổi chất lượng/độ phân giải video dựa theo băng thông mạng thời gian thực của thiết bị client.
    \item \textbf{Xây dựng bộ kiểm thử tự động:} Triển khai JUnit/Mockito kiểm thử các dịch vụ nghiệp vụ chính ở Backend, cùng Espresso để tự động kiểm thử luồng giao diện Android Client, đồng thời thiết lập tích hợp CI/CD tự động qua GitHub Actions.
    \item \textbf{Nâng cấp tính năng tương tác:} Hoàn thiện và tối ưu hóa hệ thống bình luận (bình luận lồng nhau, thời gian thực bằng WebSocket) và cho phép đánh giá xếp hạng phim (rating).
\end{itemize}

\subsection{Hướng phát triển dài hạn (>6 tháng)}
\begin{itemize}
    \item \textbf{Triển khai HTTPS và đóng gói server:} Đóng gói backend thành Docker image, vận hành sau Nginx reverse proxy với chứng chỉ Let's Encrypt; cập nhật \texttt{network\_security\_config.xml} của Android client để chỉ cho phép TLS.
    \item \textbf{CDN và bảo vệ nội dung Premium:} Tích hợp một CDN (CloudFront, Cloudflare Stream\dots) kết hợp Signed URL có thời hạn ngắn, ngăn việc chia sẻ trực tiếp liên kết video trả phí ra bên ngoài.
    \item \textbf{Push Notification:} Tích hợp Firebase Cloud Messaging (FCM) để thông báo phim mới, tập mới của series người dùng đang theo dõi và các sự kiện thể thao trực tiếp.
    \item \textbf{Hệ thống gợi ý phim:} Xây dựng module Recommendation dựa trên \texttt{WatchHistory} và \texttt{Favorite}, đề xuất nội dung phù hợp với hành vi xem của người dùng.
    \item \textbf{Mở rộng nền tảng:} Cân nhắc phát triển thêm client trên iOS hoặc một SPA web dùng chung backend, tận dụng tối đa REST API hiện có với tiền tố \texttt{/api/v1}.
\end{itemize}
